% !TEX program = xelatex
% =============================================================================
% 中文求职简历 LaTeX 模板
% Author: Kody
% Repository: https://github.com/kody1126/Chinese-resume-template-work
% Compiler: XeLaTeX
%
% 共享排版配置；正文见 tex/content，版本选择见 resume-*.tex。
% =============================================================================
\documentclass[11pt,a4paper,fontset=none]{ctexart}

% ── 字体设置 ──────────────────────────────────────────────
% 按文件路径加载仓库内固定字体；不检测或回退到操作系统字体。
\setmainfont[
    Path=fonts/,
    BoldFont=texgyretermes-bold.otf,
    ItalicFont=texgyretermes-italic.otf,
    BoldItalicFont=texgyretermes-bolditalic.otf
]{texgyretermes-regular.otf}
\setsansfont[
    Path=fonts/,
    BoldFont=texgyreheros-bold.otf,
    ItalicFont=texgyreheros-italic.otf,
    BoldItalicFont=texgyreheros-bolditalic.otf
]{texgyreheros-regular.otf}
\setmonofont[
    Path=fonts/,
    BoldFont=texgyrecursor-bold.otf,
    ItalicFont=texgyrecursor-italic.otf,
    BoldItalicFont=texgyrecursor-bolditalic.otf
]{texgyrecursor-regular.otf}
\setCJKmainfont[
    Path=fonts/,
    BoldFont=FandolHei-Regular.otf,
    ItalicFont=FandolKai-Regular.otf,
    BoldItalicFont=FandolHei-Bold.otf
]{FandolSong-Regular.otf}
\setCJKsansfont[
    Path=fonts/,
    BoldFont=FandolHei-Bold.otf,
    ItalicFont=FandolKai-Regular.otf,
    BoldItalicFont=FandolHei-Bold.otf
]{FandolHei-Regular.otf}
\setCJKmonofont[
    Path=fonts/,
    BoldFont=FandolHei-Bold.otf,
    ItalicFont=FandolKai-Regular.otf,
    BoldItalicFont=FandolHei-Bold.otf
]{FandolHei-Regular.otf}

% ── 页面布局 ──────────────────────────────────────────────
\usepackage{geometry}
\geometry{a4paper, left=1.5cm, right=1.5cm, top=1.5cm, bottom=1.5cm}

% ── 颜色定义 ──────────────────────────────────────────────
\usepackage{xcolor}
\definecolor{sectioncolor}{RGB}{25,78,121}
\definecolor{linkcolor}{RGB}{24,90,140}
\definecolor{textcolor}{RGB}{0,0,0}

% ── 标题格式 ──────────────────────────────────────────────
\usepackage{titlesec}
\titleformat{\section}
    {\Large\bfseries\color{sectioncolor}}
    {}
    {0em}
    {\uppercase}
    [\vspace{-0.8em}\rule{\textwidth}{1pt}]
\titlespacing{\section}{0pt}{0.1em}{0.1em}

% ── 列表格式 ──────────────────────────────────────────────
\usepackage{enumitem}
\setlist[itemize]{leftmargin=*, nosep, topsep=0.2em, itemsep=0.1em, parsep=0em}

% ── 图片与排版 ────────────────────────────────────────────
\usepackage{graphicx}
\setkeys{Gin}{draft=false}
\usepackage[export]{adjustbox}
\usepackage{calc}

% ── 表格 ─────────────────────────────────────────────────
\usepackage{tabularx}
\usepackage{multirow}

% ── 图标与超链接 ──────────────────────────────────────────
\usepackage{fontawesome5}
\usepackage{hyperref}
\hypersetup{colorlinks=true, linkcolor=linkcolor, urlcolor=linkcolor}

% ── 全局样式 ──────────────────────────────────────────────
\color{textcolor}
% 固定中文字距，避免不同平台的换行差异。
\xeCJKsetup{CJKglue = \hskip 0.10em}
\linespread{1.00}
\pagestyle{empty}

% ── 自定义命令 ────────────────────────────────────────────
\newcommand{\daterange}[1]{{\textbf{#1}}}
\newcommand{\entrytitle}[1]{{\textbf{#1}}}

